\paragraph{Protocol.} We study preprocessing ablations for seven-class gas identity
classification on two sensor chips (171 and 172) with common channels R1/R2.
Each variant uses the same resistance clipping, 200-point canonical temperature grid,
phase-specific validity masks, zero fill outside each output block mask, and
the same 16 statistical features per sensor-response block as the E2 tree baseline.
The reference input (fixed-weight heating and hysteresis) numerically matches
\texttt{CanonicalSubtractionFeatures}. Hyperparameters and chronological splits
are held fixed; 20 October is used for validation and 21 October for testing.
All means are over three ExtraTrees random seeds (not independent days).
Results are listed in Table~\ref{tab:e3_ablation}.

\paragraph{Results.} Relative to unreferenced log-conductance, pre-air
referencing changes test balanced accuracy by +21.2 and +26.7
percentage points (171 and 172). Adding the second air reference at equal
weight changes it by +16.8 and -0.7 points.
Compared with fixed weights, time interpolation for the heat+hysteresis
representation changes it by +2.2 and -0.0 points.
Using heat plus hysteresis rather than heat plus cooling changes it
by +2.7 and +1.8 points.

\paragraph{Interpretation.} The observed effects of air referencing and
the second air measurement are device- and protocol-dependent.
Any gain from the heat+hysteresis parameterization is a representation
effect for the evaluated ExtraTrees model, not new physical information.
Time-weighted and equal-weight air referencing are compared as alternatives
without assuming either must outperform the other.
The test day was inspected in prior exploratory analyses, so these results
are descriptive and do not establish device-independent optimality.
